\chapter{Fatiga: iniciación, vida y propagación}\label{ch:04}
\objetivos{Definir un ciclo sin ambigüedad; usar curvas S--N y modelos de tensión media con sus límites; integrar la ley de Paris y distinguir vida total de vida de propagación.}
\section{Cómo describir la carga}
La fatiga es daño asociado a cargas variables que puede llevar al fallo con tensiones máximas inferiores a la resistencia estática. Una secuencia puede contener iniciación en superficie o defecto, crecimiento de grietas pequeñas, propagación de una grieta dominante y rotura final. La proporción de vida en cada fase depende del material, superficie, tamaño y espectro. No existe una fracción universal de «vida de iniciación».

Para un ciclo entre $\sigma_{\min}$ y $\sigma_{\max}$ se definen
\begin{equation}\sigma_a=\frac{\sigma_{\max}-\sigma_{\min}}2,\quad\sigma_m=\frac{\sigma_{\max}+\sigma_{\min}}2,\quad R=\frac{\sigma_{\min}}{\sigma_{\max}}.\label{eq:ciclo}\end{equation}
Una carga completamente alternante tiene $R=-1$ y tensión media nula; una carga de cero a un máximo positivo tiene $R=0$ y $\sigma_a=\sigma_m$. La frecuencia, forma de onda, permanencias y ambiente también pueden importar. En motores aparecen ciclos de misión de baja frecuencia y vibraciones de alta frecuencia; su coexistencia exige definir qué fenómeno se modeliza.
\ilustra{ciclo-fatiga}{Definición de máximo, mínimo, amplitud y tensión media en una historia sinusoidal ilustrativa.}
\section{Enfoque tensión--vida}
Una curva S--N relaciona amplitud de tensión y número de ciclos a fallo para condiciones especificadas. Una ley de Basquin puede escribirse $\sigma_a=\sigma'_f(2N_f)^b$, donde el factor dos representa reversiones. También se emplea $\sigma_a=A N_f^b$ con otro coeficiente. No se mezclan coeficientes de ambas convenciones.
\begin{equation}N_f=\left(\frac{\sigma_a}{A}\right)^{1/b},\qquad b<0\quad\text{en el tramo ajustado}.\label{eq:basquin}\end{equation}
No se extrapola una recta de ajuste más allá de la región de datos. Algunas curvas de ciertos aceros muestran una meseta aproximada bajo condiciones concretas; no equivale a vida infinita universal. Corrosión, superficie, tamaño, entallas y temperatura pueden modificarla o eliminar su utilidad como límite. Las aleaciones de aluminio suelen caracterizarse con resistencia a una vida especificada y no con un límite universal.

La dispersión exige declarar si la curva es media o corresponde a una supervivencia determinada. Una regla empírica basada en una fracción de la resistencia a tracción puede orientar un ejercicio expresamente idealizado; no sustituye datos de fatiga de la pieza.
\section{Tensión media, entallas y deformación local}
El diagrama de Goodman modificado aproxima la penalización por tensión media de tracción:
\begin{equation}\frac{\sigma_a}{S_e}+\frac{\sigma_m}{S_u}\leq1.\label{eq:goodman}\end{equation}
$S_e$ representa aquí una resistencia alternante de referencia y $S_u$ una resistencia a tracción compatibles con el modelo. No se generaliza a compresión sin revisar el criterio, ni a plasticidad cíclica intensa. Gerber y Soderberg introducen relaciones diferentes; elegir una no es una equivalencia algebraica, sino una decisión de modelización.

La sensibilidad a entalla se resume a veces mediante $K_f=1+q(K_t-1)$, $0\leq q\leq1$ en el modelo habitual. $K_t$ es geométrico; $K_f$ se vincula al efecto de la entalla en fatiga y depende del material y del problema. No debe aplicarse dos veces un factor si la curva utilizada ya representa el detalle real.

En baja vida con plasticidad local puede emplearse la relación deformación--vida de Coffin--Manson--Basquin:
\begin{equation}\varepsilon_a=\frac{\sigma'_f}{E}(2N_f)^b+\varepsilon'_f(2N_f)^c.\label{eq:coffin}\end{equation}
Los coeficientes requieren ensayos compatibles y el cálculo puede ser implícito. Una tensión nominal elástica no describe necesariamente la amplitud de deformación en una raíz de entalla plastificada.
\section{Daño acumulado y secuencia}
La regla lineal de Miner define $D=\sum_i n_i/N_i$, siendo $N_i$ la vida que asigna el modelo al nivel $i$. El valor $D=1$ se usa como criterio idealizado de fallo. La regla no representa por sí misma efectos de orden, interacción, sobrecargas ni variación del estado superficial. En cargas complejas, el recuento de ciclos exige un procedimiento, por ejemplo rainflow, y una definición de tensión media.
\caso{M04. Un bloque no es una vida}{En un bloque ficticio hay $2000$ ciclos a un nivel con $N_1=100\,000$ y $500$ ciclos a otro con $N_2=20\,000$. Entonces $D_b=2000/100000+500/20000=0{,}045$. La estimación lineal es $1/D_b=22{,}22$ bloques. Tras 22 bloques, $D=0{,}99$; el criterio se alcanzaría durante el bloque siguiente, si se mantiene la misma idealización. El orden de los dos niveles no modifica Miner, aunque sí puede modificar el comportamiento físico.
\par\textbf{Desarrollo para comparar tu intento.} Miner suma fracciones de vida consumidas por cada nivel: $D_b=n_1/N_1+n_2/N_2$. Tras $B$ bloques completos, $D=BD_b$; el criterio ideal unitario da $B=1/D_b$.

\begin{enumerate}

\item Primera fracción de daño del bloque.

\[\fitmath{D1=\frac{2000}{100000}\approx 0.02\,}\]

\item Segunda fracción de daño.

\[\fitmath{D2=\frac{500}{20000}\approx 0.025\,}\]

\item Sumar fracciones, no vidas.

\[\fitmath{\mathrm{Db}=\frac{2000}{100000}+\frac{500}{20000}\approx 0.045\,}\]

\item Despejar bloques para daño unitario.

\[\fitmath{\mathrm{blocks}=\frac{1}{\frac{2000}{100000}+\frac{500}{20000}}\approx 22.22222\,}\]

\item Comprobar el estado tras 22 bloques completos.

\[\fitmath{D22=22\cdot \left(\frac{2000}{100000}+\frac{500}{20000}\right)\approx 0.99\,}\]

\end{enumerate}

}
\section{Propagación: ley de Paris}
En una región intermedia de crecimiento de grietas largas se utiliza
\begin{equation}\frac{da}{dN}=C(\Delta K)^m,\qquad\Delta K=Y\Delta\sigma\sqrt{\pi a}.\label{eq:paris}\end{equation}
$\Delta\sigma=\sigma_{\max}-\sigma_{\min}=2\sigma_a$. El coeficiente $C$ tiene unidades ligadas a la unidad de $a$, de $K$ y al exponente $m$. Un valor calibrado con $\mathrm{MPa\sqrt m}$ no se introduce sin conversión en una ecuación escrita con $\mathrm{Pa\sqrt m}$.

Se distinguen el umbral aparente, la región de Paris y la aceleración cerca de inestabilidad. La razón $R$, el cierre de grieta, las grietas pequeñas, las sobrecargas y el ambiente modifican el crecimiento. La ley de Paris simple no describe automáticamente esas regiones ni una fase de iniciación.

Con $Y$ y $\Delta\sigma$ constantes,
\begin{equation}N=\frac{1}{C(Y\Delta\sigma\sqrt\pi)^m}\int_{a_i}^{a_f}a^{-m/2}\,da.\label{eq:parisint}\end{equation}
Para $m\ne2$ se obtiene
\begin{equation}N=\frac{a_f^{1-m/2}-a_i^{1-m/2}}{(1-m/2)C(Y\Delta\sigma\sqrt\pi)^m};\quad m=2:\ N=\frac{\ln(a_f/a_i)}{C(Y\Delta\sigma\sqrt\pi)^2}.\label{eq:parisclosed}\end{equation}
Con $m>2$, numerador y denominador de la primera expresión son negativos y el resultado es positivo cuando $a_f>a_i$. Esta comprobación de signo evita un error frecuente. Si $Y$ cambia con $a$, se integra numéricamente la expresión completa.
\section{De la propagación al intervalo de inspección}
El tamaño final debe respetar tanto la resistencia frente a inestabilidad al máximo de carga como otros límites aplicables. Un cálculo con $\Delta K$ no sustituye la comprobación de $K_{\max}$. El tamaño inicial de inspección debe tener una base de detectabilidad y no coincidir sin más con «la resolución del equipo».

El intervalo académico puede tomarse como una fracción explícita del crecimiento calculado para explorar márgenes, pero la fracción no se presenta como un requisito reglamentario. Deben declararse fiabilidad de detección, espectro, dispersión del crecimiento y revisión de hipótesis. Los ejercicios de propagación del cuaderno separan claramente el resultado matemático de una decisión de mantenimiento.
\ilustra{sn-paris}{Curvas esquemáticas de tensión--vida y de crecimiento de grieta. Se señalan dominios de ajuste; no son datos de una aleación certificada.}

\section{Ajuste de una curva y coherencia de sus coeficientes}
Si dos puntos compatibles se ajustan con $\sigma_a=A N^b$, el logaritmo da una recta: $\ln\sigma_a=\ln A+b\ln N$. Por tanto
\begin{equation}b=\frac{\ln(\sigma_{a2}/\sigma_{a1})}{\ln(N_2/N_1)},\qquad A=\sigma_{a1}/N_1^b.\end{equation}
Los cocientes de las expresiones para $b$ son adimensionales. El coeficiente $A$ conserva la convención y unidad de tensión del ajuste. Dos puntos permiten determinar dos parámetros, pero no estimar con fiabilidad la dispersión ni demostrar que una ley de potencia sea adecuada fuera del intervalo. Los ensayos que no fallan antes de detenerse no se convierten sin más en fallos a ese número de ciclos: son observaciones censuradas.
\section{Secuencias variables y significado de una permanencia}
Rainflow convierte una historia irregular en un conjunto de ciclos según un algoritmo definido. El resultado debe conservar rangos y tensiones medias pertinentes. No se sustituye una historia por su máximo y mínimo global si contiene numerosos ciclos pequeños y grandes: se perdería la multiplicidad. Después, asignar una vida a cada clase requiere una curva compatible y un tratamiento de tensión media.

Una permanencia a alta temperatura puede introducir relajación, oxidación y daño por fluencia. Dos historias con el mismo número de máximos y mínimos pueden producir diferente daño si una mantiene la carga durante mucho más tiempo. La interacción fluencia--fatiga no queda automáticamente descrita sumando una fracción de Miner y una fracción de tiempo; cualquier regla de interacción necesita su base y límites. En un disco, las variaciones de temperatura y velocidad a lo largo de una misión también cambian el campo local de tensiones.
\section{La integral de Paris como problema de sensibilidad}
Con $m=3$ y $Y$ constante, la vida es proporcional a $1/(C\Delta\sigma^3)$. Un aumento del rango del 10\% reduce la vida del modelo al factor $1/1.1^3\approx0.7513$, manteniendo los tamaños y demás parámetros. Esta sensibilidad explica por qué confundir amplitud con rango puede introducir un error grande. La relación no debe extrapolarse si cambia el dominio de crecimiento o el tamaño final admisible.

El tamaño detectado inicial también pesa mucho: para ese exponente interviene $a_i^{-1/2}$. Reducir arbitrariamente $a_i$ prolonga matemáticamente la vida, pero puede llevar el cálculo al dominio de grietas pequeñas donde el ajuste no representa la física. La inspección y el modelo deben compartir una descripción coherente del defecto.
\ilustra{crecimiento-grieta}{Crecimiento sintético integrado con Paris. La curva acelera al crecer la grieta bajo rango de tensión constante; no incluye la fase de iniciación.}
\repaso{Una solución usa $\sigma_a$ en lugar de $\Delta\sigma$ para Paris con $m=3$. ¿Qué ocurre si todo lo demás permanece fijo?}{Como $\Delta\sigma=2\sigma_a$, la solución usa un impulsor dos veces menor y calcula una vida ocho veces mayor. Corregir la aritmética no basta: también deben comprobarse razón de carga y dominio del ajuste.}

\resumen{Un ciclo necesita máximos y mínimos. S--N estima vida bajo condiciones de ensayo; Paris estima propagación en un dominio distinto. Miner es acumulación lineal aproximada. $K_{\max}$ comprueba inestabilidad y $\Delta K$ impulsa el modelo de crecimiento; ni uno ni otro sustituye al otro.}
\fuentes{Roylance, \textit{Fatigue} y \textit{Introduction to Fracture Mechanics} \cite{MIT-FAT,MIT-FRAC}; materiales docentes de fatiga, pp. 1--37.}
